% Independently audited by root and B; source note retained.
\subsection{Strong pairs in categories with many objects}

\begin{proposition}\label{C:seven-objects-strong-pair21}
Let $C$ be a connected finite category with at least seven
objects and at most twenty-one arrows.  Assume that no
object is a retract of another.  If $C$ has a proper
strongly connected component on two objects, then
$BC\not\simeq\mathbb CP^2$.
\end{proposition}
\begin{proof}
Let $S$ be such a component, write $s$ for its number of
arrows, and put $\ell=|\operatorname{Ob}C|-2\geq5$.
Its two corner identities lie outside its minimal ideal:
an identity in that ideal would make its object a retract
of the other, by the minimal-group inverse adjustment in
Lemma~\ref{C:finite-category-Rees}.  In particular $s\geq6$.

Contract $S$ to one vertex of the undirected support graph.
The resulting graph has $\ell+1$ vertices and is connected,
so it has at least $\ell$ edges.  At least one edge meets
the contracted vertex, and strong connectivity expands
each such interface to arrows at both objects of $S$.
Counting the $\ell$ outside identities gives
\[
 |C|\geq s+2\ell+1.
\]
Thus $s\leq10$.  Every such strong pair has
$H^2(BS;\mathbb Q)=0$.  Indeed the rational categorical
Rees comparison requires both total row and total column
counts to exceed two for nonzero second homology.  Its
minimal ideal would then have at least nine arrows, in
addition to the two identities outside it, giving
$s\geq11$.

\paragraph{An extreme component.}
First suppose $S$ is a source or a sink in the
condensation order.  Take opposites to make it a source,
and let $D=C\setminus S$ have $q$ connected components.
If $H^2(BD;\mathbb Q)\ne0$, one of its connected
components, on $t$ objects, has at least $2t+2$ arrows
by Lemma~\ref{C:arrow-excess-two-H2}.  Every other
component on $t'$ objects has at least $2t'-1$ arrows.
Therefore
\[
 |D|\geq2\ell-q+3.
\]
Each component of $D$ must meet $S$, and every such
interface supplies at least two cross arrows.
Consequently
\[
 |C|\geq s+(2\ell-q+3)+2q
       =s+2\ell+q+3\geq s+2\ell+4.
\]
For $s\geq8$ and $\ell\geq5$ this exceeds twenty-one.
Thus $H^2(BD;\mathbb Q)=0$ in those cases.  Since
$H^2(BS;\mathbb Q)=0$ as well, the actual two-block
suspension comparison of
Lemma~\ref{C:two-block-zero-H2-products} makes every
degree-two product on $BC$ zero.

If $s=6$ or $7$, the minimal group is trivial and the
minimal ideal has total row and column counts both two.
There is one minimal row and one minimal column at
each object.  Lemma~\ref{C:one-row-source-block-collapse}
gives an actual homotopy-equivalent category with $S$
replaced by a point.  This removes at least $s-1$
internal arrows and at least one cross arrow.  For the
latter assertion, take an outside object receiving
arrows from $S$.  The orbit set of its right $S$-act
is an image of either nonempty stalk and therefore
has size at most the smaller stalk size; replacing
the two stalks by it saves at least one arrow.
All other arrow counts do not increase.  The new
category has at most $21-s\leq15$ arrows, contrary
to Theorem~\ref{C:CP2-lower16}.

\paragraph{A nonextreme component.}
There are distinct outside objects $a,b$ with
$a\to S\to b$.  They are distinct because a path
leaving and reentering a strong component would put
its intermediate object in that component.
Composition gives $a\to b$.  The contracted support
graph therefore contains a triangle and has at least
$\ell+1$ edges; at least two of these edges meet $S$.

Let $P_{\mathrm{ext}}$ be the sum of $|C(d,i)|-1$
and $|C(i,d)|-1$ over nonempty hom-sets with $i\in S$
and $d\notin S$.  Each support interface already
contributes two cross arrows, and these parallel
excesses are additional.  Hence
\begin{equation}\label{C:nonextreme-pair-parallel-budget}
 |C|\geq s+2\ell+3+P_{\mathrm{ext}},
 \qquad
 P_{\mathrm{ext}}\leq8-s.
\end{equation}
In particular $s\leq8$.

Suppose first $s=6$ or $7$.  The minimal ideal is the
four-arrow band with a unique arrow $k_{ij}:i\to j$
for each pair of its two objects.  Put $e_i=k_{ii}$.
It is a two-sided zero of the corner monoid at $i$,
and it differs from that corner identity.
Call the incoming condition at $i$ defective if some
cross arrow $g:d\to i$ satisfies $e_i g\ne g$.
Call the outgoing condition defective if some
$h:i\to d$ satisfies $h e_i\ne h$.
There are four conditions in total, two at each object.

An internal minimal-ideal cross arrow satisfies both
relevant conditions.  Every extra internal cross arrow
can make at most two conditions defective, one outgoing
at its source and one incoming at its target.  There
are at most $s-6$ extra internal arrows, including any
extra endomorphisms, so these contribute at most
$2(s-6)$ defects.  An external singleton hom-set is
fixed by its endpoint idempotent.  Any external defect
therefore requires a parallel hom-set and affects only
the condition at its one endpoint in $S$.
There are at most $P_{\mathrm{ext}}$ such hom-sets.
By \eqref{C:nonextreme-pair-parallel-budget}, the total
number of defective conditions is at most
\[
 2(s-6)+P_{\mathrm{ext}}\leq s-4\leq3.
\]
At least one of the four conditions is consequently
not defective.

Choose its object $i$ and the two minimal arrows to
and from the other object.  Their composite at $i$
is $e_i$.  All incoming arrows are fixed by $e_i$,
or all outgoing arrows are fixed, as appropriate.
The unit hypothesis in the one-sided deletion lemma
is automatic because $e_i$ is a zero of its corner.
Lemma~\ref{C:one-sided-corner-deletion}, or its opposite,
therefore deletes $i$ by a homotopy equivalence.
At least six arrows are removed: its corner has at
least two, there is an internal arrow in each direction,
and nonextremality supplies an external entrance and
an external exit at $i$.  The resulting category has
at most fifteen arrows, again impossible.

Finally suppose $s=8$.  Equality in the budget forces
$P_{\mathrm{ext}}=0$, so every nonempty external
hom-set is a singleton.  The minimal group is trivial,
and its band has total row and column counts either
$(2,2)$, $(3,2)$, or $(2,3)$.  These are the only
products of positive totals at least two whose band
size, plus the two outside identities, is at most
eight.  Thus one of the two total counts is two.
If it is the row count, there is one minimal row at
each object and all external incoming arrows are
fixed by the chosen local minimal idempotents.
Lemma~\ref{C:one-row-fixed-entrance-collapse} replaces
the whole block by a point, preserving homotopy type.
If it is the column count, use the dual fixed-exit
statement.  In either case the eight internal arrows
become one and no other arrow count increases.
The resulting category has at most fourteen arrows,
which is impossible.  This exhausts all proper strong
pairs under the stated hypotheses.
\end{proof}

\begin{corollary}\label{C:eight-objects-through21}
No finite category with at least eight objects and at
most twenty-one arrows has classifying space
homotopy equivalent to $\mathbb CP^2$.
\end{corollary}
\begin{proof}
Suppose such a connected category exists, and delete
retract objects until none remain.  Two deletions
would leave at most fifteen arrows and are impossible
by Theorem~\ref{C:CP2-lower16}.  Thus the remaining
category has at least seven objects.  Its excess
relative to $2k-1$ is at most six, because each
deletion lowers that excess by at least one.

If the remaining category is strongly connected,
Corollary~\ref{C:strong-category-CP2-minimum23}
excludes it.  A proper component with at least three
objects is excluded by
Proposition~\ref{C:large-SCC-minimal-candidate21}.
A proper strong pair is excluded by the preceding
proposition.  Otherwise its support is acyclic,
and Theorem~\ref{C:acyclic-arrow-excess-six} applies.
Every alternative is impossible.
\end{proof}
